\documentclass[11pt]{article}
\usepackage{preprint}

\title{The Intersection Between Race and Disability}
\author{Jiadi He\email{jh2288@georgetown.edu}}
\affil{Georgetown University}
\shorttitle{The Intersection Between Race and Disability}
\keywords{race; disability; intersectionality}

\begin{document}
\maketitle
\keywordsline

The Black experience has been captured by different people either by providing personal accounts, studying the history of the race or even oppressors rewriting the narrative to sanitize themselves. Amid all the literature there exists about Black people, one thing remains constant; struggle. For centuries, Black people have struggled to be considered human enough in the world, they have struggled for recognition, they have struggled for the bare minimum and they have struggled to ensure they provide a safe environment for their children. I cannot begin to say that I can understand the struggles of Black people and that I why I chose to view it from the lens of Black writers who have shared the Black experiences and how it has shaped them. I also seek to understand how Black people who have been oppressed for centuries have been able to deal with being disabled while Black considering that means being a member of two marginalized groups at the same time. Disability and racism are two concepts that intersect and share similar outlooks in the society and it is imperative that they are demystified so that they stop acting as obstacles.

Black people have been treated like less of human beings and oppressed for centuries and even now, they are walking around with the scars of the battles they have had to fight to get to where they are. Talila A. Lewis in the article \textit{Freeing Black Fates and Capturing Black Freedom; Reclaiming Our Humanity, Contextualizing Our Trauma and Honoring Our Resistance} discusses the inhumanity that Black people have had to endure in a world that was fashioned against them. I was greatly saddened by the statement by Lewis when she states “the terror, stress and trauma of this violence extend well beyond the individual executions and deaths. We are surrounded by reminders of the danger we are in, of the unmattering of our lives”. In the statement the author captures the fear and despair that Black people are forced to live with because of a system that despises their very existence. When I read those words, I remembered the people who died recently as a result of police brutality. The death of George Floyd and Breonna Taylor are reminders that Black people are still not free from oppression. The color of their skin is used as a weapon against them and it is leading to fear and anguish and it shows that their lives are inferior to that of others, a concept that is deeply flawed. Lewis speaks of the reality of the world we live in, a world where we have stepped back and are watching as our fellow humans are oppressed and violated because of the color of their skin. The author further explains that Black people have been taught to resist because that is the only thing that has kept them alive against the wrath of White supremacy. I found it deeply saddening that a group of people have had to live in survival mode simply because the system has been tailored against them.

As I was reading through Lewis’s text, I remembered all the biases and stereotypes that I had heard about Black people as they were depicted in movies and films. One of the things that stood out to me was the fact that the Black men were always depicted as being angry and involved in crime. The depiction works to internalize racism in the viewers since it emphasizes the connection of certain negative traits with Black people. The ways in which Black people are depicted also works to influence the way they are viewed and treated by the law. When one has been labelled a criminal even before they commit a crime, it is safe to say that their fate has been sealed. Lewis mentions that “We are not human in the wild imaginations of the “law” and nothing can change this”. She expresses the fact that Blackness has already been labelled regardless of one’s efforts to shape the trajectory of their lives and it now becomes the responsibility of Black people to protect themselves from being killed instead of asking the killers to stop killing them (Lewis). Considering the inequality that exists among Black people as a race, it is difficult to fathom how disability is handled among individuals of the Black community. The gaslighting that comes with being Black is one of the primary aspects that stands out and this makes me incline towards understanding mental disability and its intersectionality with race.

Madness and Blackness are two words that have been used together so many times as a way of erasing or discrediting the Black experience and that works against the rights of persons with disability especially in the Black community. Their Alyce Pickens in her text \textit{Black Madness::Mad Blackness} explores the ways in which disability and race intersect and how race is used as a way of discounting the Black experience while Blackness is used to discount disability (22). She explores madness not in the case of a medical condition but rather as a tool that is used for erasure and one that has been applied for ages to manage the Black community and sanitize their actions. Pickens quotes, “disability/impairment and race are neither merely biological nor wholly discursive, but rather are historical material constructs imbricated within the exploitative conditions of transnational capitalism” (20). The author explains how race and disability are defined beyond their actual meaning and how that definition is more inclined towards discrediting and disempowering Black people. Lewis in her article echoes the points raised by Pickens by pointing out “The united states has a sordid and storied history of criminalizing and deeming “insane” Black people’s natural and rational responses to the violence and terror wrought upon us by its power holders”. She points at the use of madness as a condition that is synonymous with Blackness and that works to discount both race and disability. It works to gaslit Black people and show them that their experiences do not matter.

I found both texts profoundly helpful in understanding the intersectionality of race and disability considering they both go in-depth in providing evidence that points to it. Pickens explains that disability as defined through whiteness is more universal whereas for Black people it is rather specific and shows how that affects both race and disability and the way they are understood (21). Madness is not definitive to a mental condition but is rather used as a way to provide context into Blackness and the way in which Black people react. Lewis explains that the fact that disability has been used to silence Black people and to erase their experiences for years, it has led to a internalized ableism, an aspect that has further made them more susceptible to discrimination and violence when they identified as being disabled. The articles opened my eyes to the fact that when a race is referred to as being mad/crazy because of the way they stand up for their rights and against oppression, forces a negative connotation of the word because it attacks the historical standing of the race. The result is that, disability is disregarded among people of the race and that makes them more vulnerable to the challenges posed by disability in an ableist world. Lewis advises that it is important that the Black community to get rid of the ableism mentality that has been equated to freedom and learn to embrace disability so as to better empower themselves. It is important that the gap between madness and blackness is bridged for disability discourse to be successful among Black people.

Reading both of these texts helped put racism and disability and their intersection into perspective. It allowed me to see how the oppression and violence based on race has led to a reaction by Black people that is considered madness. The labelling of experiences and reactions as madness works to negatively affect the relationship between disability and race by making Black people embrace an attitude of ableism which works against persons with disability in the community. I liked how both authors broke down the issue of race and disability and worked to champion for change.

\begin{references}[Work Cited]
\item Lewis, Talila. “Freeing Black Fates \& Capturing Black Freedom.” \textit{TALILA A. LEWIS}, 2020, \url{www.talilalewis.com/blog/freeing-black-fates-capturing-black-freedom}.
\item Pickens, Therí Alyce. \textit{Black Madness:: Mad Blackness}. Duke University Press, 2019.
\end{references}

\end{document}
